\Needspace{8\baselineskip}
\Article{264}
Auf diese Weise werden also für jede gegebene Nichtquadratzahl als Determinante alle zuweisbaren Charaktere gleichmäßig in zwei Arten $P$, $Q$ verteilt: Keinem der Charaktere $Q$ kann eine eigentlich primitive (positive) Form entsprechen, während bei den übrigen $P$, soweit wir bisher wissen, nichts verhindert, dass sie zu solchen Formen gehören. Bei diesen beiden Arten beachte man besonders die folgende Aussage, die leicht aus ihrem Kriterium folgt: Setzt man einen Charakter aus $P$ mit einem Charakter aus $Q$ zusammen (nach der Regel von Artikel 246, so als entspräche auch diesem ein Geschlecht), so entsteht ein Charakter aus $Q$; setzt man dagegen zwei Charaktere aus $P$. oder zwei aus $Q$ zusammen, so gehört der entstehende Charakter zu $P$. Mithilfe dieses Satzes lässt sich auch bei negativen und uneigentlich primitiven Geschlechtern die Hälfte aller zuweisbaren Charaktere folgendermaßen ausschließen.

I. Bei negativer Determinante $D$ sind die negativen Geschlechter in dieser Hinsicht den positiven völlig entgegengesetzt: Keiner der Charaktere $P$ gehört zu einem negativen eigentlich primitiven Geschlecht, sondern alle diese Geschlechter haben Charaktere aus $Q$. Gilt nämlich $D'\equiv1\ (\mathrm{mod.}\ 4)$, so ist $-D'$ eine positive Zahl der Form $4n+3$; somit gibt es unter $a$, $b$, $c$ usw. eine ungerade Anzahl von Zahlen der Form $4n+3$, von denen jede $-1$ als Nichtrest hat. Daher enthält der Gesamtcharakter der Form $(-1,0,D)$ in diesem Fall eine ungerade Anzahl von Einzelcharakteren aus $\Omega$, gehört also zu $Q$. Gilt $D'\equiv3\ (\mathrm{mod.}\ 4)$, so findet sich aus demselben Grund unter $a$, $b$, $c$ usw. entweder keine Zahl der Form $4n+3$ oder zwei oder vier usw. Da aber in diesem Fall entweder $3,4$ oder $3,8$ oder $7,8$ unter den Einzelcharakteren der Form $(-1,0,D)$ vorkommt, gehört der Gesamtcharakter dieser Form auch hier\textperiodcentered{}zu $Q$. Dasselbe folgt ebenso leicht in den übrigen Fällen
\SourceLeaf{302}
sodass die negative Form $(-1, 0, D)$ stets einen Charakter aus $Q$ hat. Da diese Form mit jeder anderen negativen eigentlich primitiven Form derselben Determinante zusammengesetzt eine solche positive Form ergibt, kann offenbar keine negative eigentlich primitive Form einen Charakter aus $P$ haben.

II. Für uneigentlich primitive (positive) Geschlechter beweist man ebenso, dass es sich entweder genauso verhält wie bei den eigentlich primitiven oder entgegengesetzt, je nachdem, ob $D\equiv1$ oder $\equiv5\;(\mathrm{mod.}\,8)$ gilt. Im ersten Fall gilt nämlich auch $D'\equiv1\;(\mathrm{mod.}\,8)$; daher gibt es unter den Zahlen $a$, $b$, $c$ usw. entweder keine Zahl der Formen $8n+3$ und $8n+5$ oder zwei oder vier usw. (denn ein Produkt aus beliebig vielen ungeraden Zahlen, unter denen die Zahlen der Formen $8n+3$ und $8n+5$ zusammen in ungerader Anzahl vorkommen, wird stets entweder $\equiv3$ oder $\equiv5\;(\mathrm{mod.}\,8)$, während das Produkt aller $a$, $b$, $c$ usw. entweder $D'$ oder $-D'$ gleich sein muss). Daher enthält der Gesamtcharakter der Form $(2, 1, \frac{1-D}{2})$ entweder keinen Einzelcharakter aus $\Omega$ oder zwei oder vier usw., gehört also zu $P$. Da nun jede uneigentlich primitive (positive) Form der Determinante $D$ als zusammengesetzt aus $(2, 1, \frac{1-D}{2})$ und einer eigentlich primitiven (positiven) Form derselben Determinante angesehen werden kann, hat in diesem Fall offenbar keine uneigentlich primitive (positive) Form einen Charakter aus $Q$. Im anderen Fall. $D\equiv5\;(\mathrm{mod.}\,8)$ ist alles entgegengesetzt: $D'$, das ebenfalls $\equiv5$ ist, enthält gewiss eine ungerade Anzahl von Faktoren der Formen $8n+3$ und $8n+5$. Daher gehört der Charakter der Form $(2, 1, \frac{1-D}{2})$ und somit auch der Charakter jeder uneigentlich primitiven (positiven) Form der Determinante $D$ zu $Q$; also kann keinem der Charaktere $P$ ein uneigentlich primitives positives Geschlecht entsprechen.

III. Bei negativer Determinante sind schließlich die negativen uneigentlich primitiven Geschlechter wiederum den positiven uneigentlich primitiven entgegengesetzt: Die ersten können keinen Charakter aus $P$ oder aus $Q$ haben, je nachdem, ob $D\equiv1$ oder $\equiv5\;(\mathrm{mod.}\,8)$ gilt, also je nachdem, ob $-D$ von der Form $8n+7$ oder $8n+3$ ist. Denn die Komposition der Form $(-1, 0, D)$, deren Charakter aus $Q$ ist, mit negativen uneigentlich primitiven Formen derselben Determinante ergibt positive uneigentlich primitive Formen. Sind also bei diesen die Charaktere $Q$ ausgeschlossen, so müssen bei jenen die Charaktere $P$ ausgeschlossen sein und umgekehrt.

\Needspace{11\baselineskip}
\SourceLeaf{303}
\begin{center}
\textit{Eine besondere Methode, Primzahlen in zwei Quadrate zu zerlegen.}
\end{center}

\Needspace{8\baselineskip}
\Article{265}
Aus den Untersuchungen der Artikel 257, 258 über die Anzahl der ambigen Klassen, auf denen alles Vorige beruht, lassen sich viele weitere sehr beachtenswerte Folgerungen ziehen, die wir der Kürze wegen weglassen müssen. Die folgende ist jedoch so elegant, dass wir sie nicht übergehen können. Für eine positive Determinante $p$, die eine Primzahl der Form $4n+1$ ist, haben wir gezeigt, dass es nur eine eigentlich primitive ambige Klasse gibt; daher sind alle eigentlich primitiven ambigen Formen einer solchen Determinante eigentlich äquivalent. Ist also $b$ die unmittelbar unter $\surd p$ liegende positive ganze Zahl und $p-bb
=a'$, so sind die Formen $(1, b, -a')$, $(-1, b, a')$ eigentlich äquivalent; da beide offenbar reduziert sind, liegt die eine in der Periode der anderen. Geben wir der ersten Form in ihrer Periode den Index $0$, so muss der Index der zweiten ungerade sein (da die ersten Koeffizienten beider Formen entgegengesetzte Vorzeichen haben); er sei also $=2m+1$. Sind ferner die Formen mit den Indizes $1, 2, 3$ usw. jeweils
\[
(-a', b', a''),\quad (a'', b'', -a'''),\quad (-a''', b''', a'''')\text{ usw.}
\]
so entsprechen offenbar den Indizes $2m$, $2m-1$, $2m-2$, $2m-3$ usw. die Formen
\[
(a', b, -1),\quad (-a'', b', a'),\quad (a''', b'', -a''),\quad (-a'''', b''', a''')\text{ usw.}
\]
Ist also die Form mit dem Index $m$ gleich $(A, B, C)$, so ist sie auch $(-C, B, -A)$; daher gilt $C=-A$ und $p=BB+AA$. Somit lässt sich jede Primzahl der Form $4n+1$ in zwei Quadrate zerlegen (eine Aussage, die wir oben in Artikel 182 aus völlig anderen Prinzipien hergeleitet haben). Zu einer solchen Zerlegung gelangt man durch eine sehr einfache und ganz einheitliche Methode: Man entwickelt die Periode der reduzierten Form, deren Determinante jene Primzahl und deren erster Koeffizient $1$ ist, bis zu einer Form, deren äußere Koeffizienten dem Betrag nach gleich und im Vorzeichen entgegengesetzt sind. So erhält man \textit{zum Beispiel}\ für $p=233$ $(1, 15, -8)$, $(-8, 9, 19)$, $(19, 10, -7)$, $(-7, 11, 16)$, $(16, 5, -13)$, $(-13, 8, 13)$ und $233=64+169$. Ferner muss $A$ offenbar ungerade sein (da $(A, B, -A)$ eigentlich primitiv sein muss), also $B$ gerade. --- Da bei einer positiven Determinante $p$, die eine Primzahl der Form $4n+1$ ist, auch die uneigentlich primitive Ordnung nur eine ambige Klasse enthält, sind offenbar, wenn $g$ die
\SourceLeaf{304}
unmittelbar unter $\surd p$ liegende ungerade Zahl und $p-gg=4h$ ist, die reduzierten uneigentlich primitiven Formen $(2, g, -2h)$, $(-2, g, 2h)$ eigentlich äquivalent; somit liegt die eine in der Periode der anderen. Durch ganz ähnliche Überlegungen wie zuvor folgt, dass die Periode der Form $(2, g, -2h)$ eine Form enthält, deren äußere Koeffizienten dem Betrag nach gleich und im Vorzeichen entgegengesetzt sind, sodass sich auch daraus die Zerlegung der Zahl $p$ in zwei Quadrate gewinnen lässt. Die äußeren Koeffizienten dieser Form sind offenbar gerade, der mittlere also ungerade. Da eine Primzahl nur auf eine Weise in zwei Quadrate zerlegt werden kann, ist die nach dieser zweiten Methode gefundene Form entweder $(B, \pm A, -B)$ oder $(-B, \pm A, B)$. So erhält man in unserem Beispiel für $p=233$ $(2, 15, -4)$, $(-4, 13, 16)$, $(16, 3, -14)$, $(-14, 11, 8)$, $(8, 13, -8)$ und $233=169+64$ wie oben.
